\honest{The Power of Intelligence}{Eliezer Yudkowsky}{2007}

In our skulls we carry three pounds of wet grey tissue. It does not look dangerous. Aristotle thought the brain cooled the blood. \nb{Roughly so: in \textsc{Parts of Animals} the brain ``tempers the heat and seething of the heart.''}

Five million years ago predators and prey fought with teeth, claws, shells and venom. Every weapon met a defense, and no species ruled more than its own niche. Then came the Squishy Things, with no armor, no claws and weak fingers. Had you watched evolution at work, you would have thought it absurd that they could cut stone, melt metal, set off a nuclear explosion or handle DNA. To evolution a year is nothing, a mere six hundred trillion trillion trillion trillion Planck intervals, and by that reckoning the Squishy Things would not fly for ten million years.

``Now explain to me again why an Artificial Intelligence can't do anything interesting over the Internet unless a human programmer builds it a robot body.'' \nb{Every feat in the parable began with hands, however weak. The AI in the question has no body, and the post does not say how it would act on the world.}

People who hear ``intelligence,'' I have observed, often think of book smarts. So they say, ``It takes more than intelligence to succeed professionally,'' as if charisma resided in the kidneys. \nb{The sayings use ``intelligence'' in the narrow sense; the post's reply counts everything a brain does, charisma included. Whether an AI built for the narrow kind would have the rest is what the sayings question, and the wide definition settles it without argument.} They say intelligence is no match for a gun, as if guns grew on trees, and ask where an AI would get money, as if the first humans found dollar bills in the forest. The human species imagined money into existence. \nb{That the species made guns and money, over long periods, does not say what one new intelligence could do where others already have them.} The archetype of intelligence, I keep explaining, is not the savant of \textsc{Rain Man} but a human being.

Venture capitalists ask how the Machine Intelligence Research Institute would commercialize a true AI. \nb{The institute took that name in 2013; the sentence has been revised since the post was first published in 2007.} That is a framing problem. People can imagine each feat of intelligence, but not one power that could do them all. Promise a visit to Mars, a proof of the Riemann Hypothesis and cures for cancer, aging and stupidity, and it ``just sounds wrong.'' I say it sounds wrong because the list is too short: ``Who could have imagined, ever so long ago, what minds would someday do?''

One trick cured smallpox, built airplanes and tamed fire. Our ignorance of how it works is a fact about us: ``A blank map does not correspond to a blank territory.'' Intelligence is as real as electricity, only far more powerful, far more dangerous, and ``a tiny little bit harder'' to build a generator for.
